% ==========================================================================
% examen-espana.tex — Examen de ejemplo del dictamen de coyuntura: España
% Guía del dictamen (1D01), entregable 3. Autor: Víctor Gutiérrez Marcos.
% Formato de la convocatoria de 2026 (1 h 45 min, cinco cuestiones, cuadernillo
% de cuadros y gráficos). Datos reales y oficiales al 26/10/2026 (datos/).
% Cuadros: cuadros/*.tex, generados con generar_cuadros.py desde datos/*.csv.
% Gráficos: graficos/*.tex (pgfplots, paleta «dic» de latex/tikz-tcee.tex).
% Compilación: ver README.md.
% ==========================================================================
\documentclass[11pt,a4paper]{article}
\usepackage[table]{xcolor}
\input{tikz-tcee}                        % fuentes (newpx), colores dic…, capas
\usepackage[utf8]{inputenc}
\usepackage[spanish,es-noshorthands,es-nodecimaldot]{babel}
\usepackage[a4paper,top=1.9cm,bottom=1.8cm,left=2.2cm,right=2.2cm,headheight=14pt]{geometry}
\usepackage{pgfplotstable}
\usepackage{booktabs,array}
\usepackage{enumitem}
\usepackage{fancyhdr}
\usepackage{lastpage}
\usepackage{microtype}
\usepackage[hidelinks,pdfusetitle]{hyperref}
\hypersetup{pdftitle={Dictamen de coyuntura: España (examen de ejemplo)},pdfauthor={Víctor Gutiérrez Marcos}}

% --- Cabecera y pie, sobrios, como el cuadernillo oficial ---------------------
\pagestyle{fancy}
\fancyhf{}
\renewcommand{\headrulewidth}{0pt}
\fancyhead[L]{\footnotesize\textcolor{tceegris}{Oposición TCEE · Primer ejercicio, segunda parte: dictamen}}
\fancyhead[R]{\footnotesize\textcolor{tceegris}{Examen de ejemplo (no oficial)}}
\fancyfoot[C]{\footnotesize Página \thepage\ de \pageref{LastPage}}

% --- Colores añadidos para este cuadernillo (CRITERIOS §7: se definen una vez) --
% Balanza de pagos, gama fría del sector exterior: bienes (azul medio), viajes
% (azul claro) y resto de servicios (azul claro rayado); rentas (gris); saldos
% (negro). PIIN y deuda externa (marrón, sin otro uso). Cuentas públicas (pizarra).
\definecolor{dicbienes}{HTML}{2C6E9B}
\definecolor{dicturismo}{HTML}{8CC3E6}
\definecolor{dicrentas}{HTML}{A6A6A6}
\definecolor{dicsaldo}{HTML}{1A1A1A}
\definecolor{dicpiin}{HTML}{8C510A}
\definecolor{dicfiscal}{HTML}{4A5A6A}
\definecolor{cabecera}{HTML}{E8E8E8}
\definecolor{fondoprev}{HTML}{F4F4F4}

\tikzset{dicbarrarayada/.style={fill=#1, draw=#1!65!black, line width=0.3pt,
    postaction={pattern=north east lines, pattern color=dicbienes}}}
\pgfplotsset{
  dicgraf/.style={width=12.0cm, height=\altograf, scale only axis=false,
    axis lines*=left, tick align=outside, clip=false,
    ymajorgrids, grid style={tceegris!25},
    ticklabel style={font=\scriptsize}, label style={font=\scriptsize},
    xticklabel style={/pgf/number format/1000 sep=},
    yticklabel style={/pgf/number format/.cd, use comma, 1000 sep={}},
    legend style={at={(xticklabel cs:0.5)}, anchor=north, yshift=-3pt, draw=none, font=\scriptsize,
      /tikz/every even column/.append style={column sep=7pt}},
    legend cell align=left},
  leyendalinea/.style={legend image code/.code={\draw[##1] (0cm,0cm) -- (0.55cm,0cm);}},
}

% --- Cuadros -------------------------------------------------------------------
\newlength{\altograf}\setlength{\altograf}{5.6cm}
\newlength{\anchonum}\newlength{\anchoetq}\setlength{\anchoetq}{6.3cm}
\newcolumntype{N}{>{\raggedleft\arraybackslash}p{\anchonum}}
\newcommand{\cuadro}[3]{\par\noindent{\small\textbf{Cuadro #1.} #2 \textit{(#3)}}\par\vspace{2pt}}
\newcommand{\fuentecuadro}[1]{\par\vspace{1pt}\noindent{\tiny\textit{Fuente:} #1}\par}
\newcommand{\cabprev}[1]{\cellcolor{cabecera}\textbf{#1\,P}}
\newcommand{\prevcel}[1]{\cellcolor{fondoprev}#1}
\newenvironment{bloquecuadro}{\par\noindent\begingroup\scriptsize\setlength{\tabcolsep}{1.6pt}\renewcommand{\arraystretch}{1.08}}{\endgroup\par}

% --- Gráficos: \grafico{n}{título}{fuente}{fichero} ---------------------------------
\newcommand{\grafico}[4]{%
  \begin{minipage}[t]{0.495\linewidth}\centering
    {\small\textbf{Gráfico #1.} #2\par}\vspace{1pt}
    \input{graficos/#4}\par\vspace{-2pt}
    \begin{minipage}{0.95\linewidth}\raggedright\tiny\textit{Fuente:} #3\end{minipage}
  \end{minipage}}

\setlist{itemsep=3pt,topsep=4pt}
\setlength{\parindent}{0pt}

\begin{document}

% ==========================================================================
% Página 1: instrucciones
% ==========================================================================
\thispagestyle{fancy}
\begin{center}
{\large\textbf{Oposición al Cuerpo Superior de Técnicos Comerciales y Economistas del Estado}}\\[4pt]
{\normalsize Primer ejercicio, segunda parte: dictamen sobre un escenario macroeconómico}\\[10pt]
{\Large\textbf{Examen de ejemplo: la economía española}}\\[4pt]
{\small Preparado por Víctor Gutiérrez Marcos con el formato de la convocatoria de 2026 · No es un examen oficial}\\[2pt]
{\small Datos oficiales disponibles el 26 de octubre de 2026}
\end{center}
\vspace{8pt}
\rule{\linewidth}{0.4pt}
\vspace{2pt}

\textbf{Instrucciones}

\begin{itemize}[leftmargin=1.2em]
\item El ejercicio consta de \textbf{cinco cuestiones}. El tribunal valorará fundamentalmente \textbf{las capacidades analítica y de síntesis del aspirante, así como su claridad expositiva}.
\item Dispone de \textbf{1 hora y 45 minutos} para responder a las cuestiones planteadas. No se podrá dar la vuelta al cuadernillo ni acceder a su contenido hasta que lo indique el tribunal; el incumplimiento de este punto supondrá la eliminación del aspirante.
\item Las respuestas se redactarán en las hojas autocopiativas depositadas sobre la mesa, con bolígrafo de tinta azul o negra. En la parte superior de cada hoja deberán consignarse el nombre y los apellidos antes de que finalice la prueba.
\item No escriba con abreviaturas, ya que la lectura se realizará de forma literal. Solo se admiten las abreviaturas que figuran en el propio enunciado y las convencionales del análisis económico (por ejemplo, PIB, pp, pb o FMI).
\item \textbf{No está permitido el uso de calculadora.} Está permitido el uso de tapones para los oídos desde el comienzo de la prueba, pero no el de cascos o auriculares de ningún tipo. Los dispositivos electrónicos (móviles, relojes inteligentes, etc.) deberán permanecer apagados y fuera del alcance de los aspirantes.
\item El tribunal no responderá a preguntas, para garantizar la igualdad de oportunidades. En caso de duda, cada aspirante debe interpretar lo que está escrito.
\item No será posible abandonar la sala durante los 10 primeros minutos ni durante los 10 últimos minutos de la prueba.
\item Cuando se anuncie el final de la prueba, los aspirantes deberán dejar de escribir inmediatamente; el incumplimiento de este punto supondrá la eliminación del aspirante.
\end{itemize}

\vspace{6pt}
\textbf{Contenido del cuadernillo}

\begin{itemize}[leftmargin=1.2em]
\item Enunciado (página 2).
\item Cuadros 1 a 7 y cuadro 2 bis: datos anuales de 2007 a 2025 y, en algunas filas, previsiones de la Comisión Europea para 2026 y 2027 (columnas marcadas con P y sombreadas).
\item Gráficos 1 a 12: datos anuales, trimestrales y mensuales; los más recientes, hasta el tercer trimestre de 2026.
\item Salvo que se indique otra cosa, «zona del euro» se refiere a los veinte países que la formaban en 2025 y «UE», a los veintisiete Estados miembros. Las celdas en blanco corresponden a datos no disponibles.
\end{itemize}

\vfill
\begin{center}
\textbf{NO DÉ LA VUELTA A LAS HOJAS HASTA QUE SE LO INDIQUE EL TRIBUNAL}
\end{center}
\vspace{10pt}

\clearpage
% ==========================================================================
% Página 2: enunciado
% ==========================================================================
En los cuadros y gráficos a continuación de este enunciado se facilita información económica sobre \textbf{España} referida a los años naturales 2007 a 2025, con previsiones de la Comisión Europea para 2026 y 2027 y datos trimestrales y mensuales hasta 2026. Atendiendo a esa información, elabore un escrito respondiendo a las siguientes cuestiones:

\begin{enumerate}[leftmargin=1.6em, label=\textbf{\arabic*.}, itemsep=9pt]
\item Describa brevemente el \textbf{perfil cíclico} de la economía española a lo largo del periodo 2007-2025, desde la óptica de la demanda y de la oferta, y valore su situación coyuntural a la luz de los datos trimestrales de 2026 y de las previsiones para 2026 y 2027. ¿Qué factores podrían explicar el dinamismo de la inversión en los últimos años y su distinta composición respecto de la fase expansiva anterior a 2008? ¿A qué podría deberse que, desde 2019, el PIB por habitante haya crecido bastante menos que el PIB?

\item Caracterice el \textbf{mercado laboral} español a lo largo de la serie. ¿Puede afirmarse que se encuentra tensionado? ¿A qué podría deberse la caída de la temporalidad desde 2022? Analice la evolución de la \textbf{inflación} desde 2021 y sus principales determinantes. ¿Qué factores podrían explicar que en 2026 la inflación española se sitúe por encima de la de la zona del euro? ¿Aprecia riesgo de efectos de segunda vuelta?

\item Analice la evolución del \textbf{sector exterior} a lo largo de todo el horizonte temporal, identificando sus fortalezas y debilidades. ¿Considera que la corrección del desequilibrio exterior posterior a 2008 tiene carácter estructural o cíclico? ¿Por qué la posición de inversión internacional neta empeora en 2025 pese al superávit exterior? ¿Cree que la política arancelaria de la actual administración de Estados Unidos ha tenido algún impacto sobre la economía española?

\item Valore la orientación de la \textbf{política monetaria} y su transmisión a la economía española, en particular la reacción del Banco Central Europeo ante el repunte de la inflación de 2026. Valore brevemente la situación del \textbf{sistema bancario}. Evalúe asimismo el tono de la \textbf{política fiscal}, su relación con el ciclo y su combinación con la política monetaria. ¿Considera que existen riesgos apreciables para la sostenibilidad de la deuda pública en el marco de las reglas fiscales europeas vigentes?

\item Identifique los principales \textbf{desequilibrios y vulnerabilidades} de la economía española y los riesgos al alza y a la baja que rodean las previsiones para 2026 y 2027. Teniendo en cuenta las tendencias demográficas, ¿qué \textbf{reformas estructurales y medidas de política económica} recomendaría para sostener el crecimiento a medio y largo plazo?
\end{enumerate}

\clearpage
% ==========================================================================
% Cuadernillo: páginas apaisadas
% ==========================================================================
% (A4 apaisado desde aquí: se cambia el tamaño de la página y la caja de texto)
\paperwidth=297mm \paperheight=210mm \pdfpagewidth=297mm \pdfpageheight=210mm
\newgeometry{layoutwidth=297mm,layoutheight=210mm,top=1.5cm,bottom=1.4cm,left=1.4cm,right=1.4cm,headheight=14pt,headsep=8pt,footskip=16pt}
\fancyheadoffset{0pt}

% ---- Página 3: cuadros 1 y 2 --------------------------------------------------
\begin{bloquecuadro}% Generado por generar_cuadros.py a partir de datos/cuadro1.csv; no editar a mano.
\cuadro{1}{Demanda agregada y actividad}{variación anual en volumen, en \%, salvo indicación}
\setlength{\anchonum}{\dimexpr(\linewidth-\anchoetq)/21-2\tabcolsep\relax}
\begin{tabular}{@{}>{\raggedright\arraybackslash}p{\dimexpr\anchoetq-\tabcolsep\relax}*{21}{N}@{}}
\toprule
\rowcolor{cabecera} & \textbf{2007} & \textbf{2008} & \textbf{2009} & \textbf{2010} & \textbf{2011} & \textbf{2012} & \textbf{2013} & \textbf{2014} & \textbf{2015} & \textbf{2016} & \textbf{2017} & \textbf{2018} & \textbf{2019} & \textbf{2020} & \textbf{2021} & \textbf{2022} & \textbf{2023} & \textbf{2024} & \textbf{2025} & \cabprev{2026} & \cabprev{2027} \\
\midrule
\textbf{PIB} & 3{,}5 & 0{,}8 & $-$3{,}8 & 0{,}1 & $-$0{,}6 & $-$2{,}9 & $-$1{,}4 & 1{,}5 & 4{,}1 & 2{,}9 & 2{,}9 & 2{,}4 & 2{,}0 & $-$10{,}9 & 6{,}7 & 6{,}4 & 2{,}4 & 3{,}7 & 2{,}6 & \prevcel{2{,}4} & \prevcel{1{,}9} \\
\hspace*{0.8em}Consumo privado & 3{,}5 & $-$0{,}6 & $-$3{,}5 & 0{,}5 & $-$2{,}4 & $-$3{,}3 & $-$2{,}9 & 1{,}7 & 3{,}0 & 2{,}7 & 3{,}1 & 1{,}7 & 1{,}1 & $-$12{,}1 & 7{,}1 & 4{,}9 & 1{,}9 & 3{,}1 & 3{,}1 & \prevcel{2{,}9} & \prevcel{2{,}0} \\
\hspace*{0.8em}Consumo público & 5{,}1 & 4{,}9 & 3{,}8 & 0{,}7 & 0{,}6 & $-$3{,}9 & $-$2{,}3 & $-$0{,}5 & 1{,}9 & 0{,}9 & 1{,}0 & 2{,}1 & 2{,}2 & 3{,}5 & 3{,}6 & 0{,}8 & 4{,}3 & 3{,}3 & 2{,}4 & \prevcel{1{,}7} & \prevcel{1{,}8} \\
\hspace*{0.8em}Formación bruta de capital fijo & 4{,}0 & $-$4{,}2 & $-$17{,}1 & $-$5{,}1 & $-$7{,}3 & $-$7{,}1 & $-$3{,}5 & 4{,}4 & 5{,}3 & 2{,}0 & 6{,}8 & 6{,}5 & 4{,}9 & $-$8{,}9 & 2{,}6 & 4{,}2 & 5{,}3 & 4{,}9 & 5{,}4 & \prevcel{4{,}4} & \prevcel{2{,}8} \\
\hspace*{1.6em}\textit{Viviendas} & 0{,}7 & $-$8{,}5 & $-$20{,}2 & $-$13{,}1 & $-$12{,}4 & $-$5{,}5 & $-$7{,}8 & 9{,}6 & $-$3{,}4 & 8{,}8 & 14{,}7 & 14{,}5 & 7{,}0 & $-$8{,}6 & 0{,}8 & 9{,}0 & 0{,}3 & 1{,}6 & 3{,}6 &  &  \\
\hspace*{1.6em}\textit{Otras construcciones} & 3{,}1 & $-$2{,}3 & $-$11{,}4 & $-$8{,}6 & $-$11{,}7 & $-$14{,}2 & $-$7{,}6 & $-$1{,}6 & 7{,}2 & $-$5{,}4 & $-$0{,}7 & 5{,}3 & 10{,}0 & $-$8{,}3 & 0{,}1 & $-$1{,}9 & 9{,}0 & 9{,}9 & 6{,}5 &  &  \\
\hspace*{1.6em}\textit{Maquinaria, bienes de equipo y armamento} & 8{,}4 & $-$3{,}4 & $-$26{,}1 & 4{,}7 & $-$0{,}7 & $-$6{,}7 & 2{,}9 & 6{,}2 & 9{,}3 & 1{,}6 & 9{,}1 & 4{,}2 & 1{,}7 & $-$13{,}5 & 3{,}3 & 2{,}1 & 3{,}0 & 3{,}0 & 6{,}1 &  &  \\
\hspace*{1.6em}\textit{Productos de la propiedad intelectual} & 8{,}7 & 5{,}7 & $-$1{,}8 & 10{,}0 & 2{,}0 & 3{,}8 & $-$1{,}0 & 4{,}4 & 5{,}6 & 5{,}7 & 2{,}7 & 1{,}3 & 0{,}4 & $-$2{,}7 & 8{,}4 & 9{,}4 & 12{,}6 & 7{,}1 & 5{,}9 &  &  \\
\hspace*{0.8em}Exportaciones de bienes y servicios & 7{,}7 & $-$0{,}9 & $-$10{,}8 & 9{,}0 & 8{,}2 & 0{,}9 & 4{,}3 & 4{,}5 & 4{,}4 & 5{,}4 & 5{,}6 & 1{,}7 & 2{,}3 & $-$20{,}1 & 13{,}4 & 14{,}2 & 2{,}2 & 3{,}8 & 2{,}6 & \prevcel{1{,}1} & \prevcel{2{,}3} \\
\hspace*{0.8em}Importaciones de bienes y servicios & 8{,}2 & $-$5{,}5 & $-$18{,}3 & 6{,}2 & $-$0{,}6 & $-$5{,}9 & $-$0{,}3 & 6{,}8 & 5{,}1 & 2{,}6 & 6{,}7 & 3{,}9 & 1{,}3 & $-$15{,}1 & 15{,}0 & 7{,}7 & $-$0{,}1 & 3{,}3 & 5{,}5 & \prevcel{2{,}4} & \prevcel{2{,}8} \\
Demanda nacional (contribución, pp) & 4{,}2 & $-$0{,}8 & $-$6{,}5 & $-$0{,}5 & $-$2{,}9 & $-$4{,}9 & $-$2{,}9 & 2{,}0 & 4{,}2 & 1{,}9 & 3{,}0 & 3{,}0 & 1{,}6 & $-$8{,}8 & 6{,}9 & 4{,}1 & 1{,}4 & 3{,}3 & 3{,}4 &  &  \\
\hspace*{0.8em}de la cual, variación de existencias (pp) & 0{,}1 & 0{,}0 & $-$0{,}5 & 0{,}3 & 0{,}0 & $-$0{,}7 & 0{,}0 & 0{,}3 & 1{,}0 & $-$0{,}2 & $-$0{,}2 & 0{,}4 & $-$0{,}4 & $-$0{,}7 & 1{,}6 & 0{,}3 & $-$1{,}5 & 0{,}0 & 0{,}2 &  &  \\
Demanda externa neta (contribución, pp) & $-$0{,}6 & 1{,}5 & 2{,}8 & 0{,}6 & 2{,}3 & 2{,}0 & 1{,}4 & $-$0{,}5 & $-$0{,}1 & 1{,}0 & $-$0{,}1 & $-$0{,}6 & 0{,}4 & $-$2{,}2 & $-$0{,}3 & 2{,}3 & 0{,}9 & 0{,}3 & $-$0{,}8 &  &  \\
PIB por habitante & 1{,}5 & $-$0{,}9 & $-$4{,}6 & $-$0{,}3 & $-$1{,}0 & $-$2{,}9 & $-$1{,}1 & 1{,}8 & 4{,}2 & 2{,}8 & 2{,}7 & 1{,}9 & 1{,}2 & $-$11{,}4 & 6{,}7 & 5{,}4 & 1{,}2 & 2{,}6 & 1{,}6 &  &  \\
PIB nominal (miles de millones de euros) & 1078 & 1112 & 1073 & 1077 & 1069 & 1036 & 1026 & 1039 & 1087 & 1123 & 1170 & 1212 & 1254 & 1129 & 1235 & 1376 & 1496 & 1599 & 1690 & \prevcel{1773} & \prevcel{1853} \\
PIB potencial & 3{,}8 & 3{,}0 & 1{,}2 & 1{,}3 & 0{,}7 & $-$0{,}3 & $-$0{,}5 & $-$0{,}1 & 0{,}1 & 0{,}3 & 0{,}5 & 0{,}7 & 1{,}0 & 0{,}3 & 1{,}0 & 1{,}9 & 2{,}4 & 2{,}5 & 2{,}7 & \prevcel{2{,}7} & \prevcel{2{,}5} \\
Brecha de producción (\% del PIB potencial) & 3{,}3 & 1{,}1 & $-$3{,}9 & $-$5{,}1 & $-$6{,}3 & $-$8{,}7 & $-$9{,}6 & $-$8{,}1 & $-$4{,}5 & $-$1{,}9 & 0{,}4 & 2{,}1 & 3{,}1 & $-$8{,}5 & $-$3{,}3 & 0{,}9 & 1{,}0 & 1{,}9 & 2{,}0 & \prevcel{1{,}7} & \prevcel{1{,}2} \\
\bottomrule
\end{tabular}
\fuentecuadro{Eurostat (contabilidad nacional del INE); PIB potencial, brecha de producción y previsiones: Comisión Europea (AMECO, previsiones de primavera de 2026). P: previsión.}
\end{bloquecuadro}
\vspace{10pt}
\begin{bloquecuadro}% Generado por generar_cuadros.py a partir de datos/cuadro2.csv; no editar a mano.
\cuadro{2}{Oferta: valor añadido bruto por ramas de actividad}{variación anual en volumen, en \%}
\setlength{\anchonum}{\dimexpr(\linewidth-\anchoetq)/19-2\tabcolsep\relax}
\begin{tabular}{@{}>{\raggedright\arraybackslash}p{\dimexpr\anchoetq-\tabcolsep\relax}*{19}{N}@{}}
\toprule
\rowcolor{cabecera} & \textbf{2007} & \textbf{2008} & \textbf{2009} & \textbf{2010} & \textbf{2011} & \textbf{2012} & \textbf{2013} & \textbf{2014} & \textbf{2015} & \textbf{2016} & \textbf{2017} & \textbf{2018} & \textbf{2019} & \textbf{2020} & \textbf{2021} & \textbf{2022} & \textbf{2023} & \textbf{2024} & \textbf{2025} \\
\midrule
\textbf{Valor añadido bruto total} & 4{,}0 & 1{,}1 & $-$3{,}3 & $-$0{,}2 & $-$0{,}3 & $-$2{,}8 & $-$1{,}3 & 1{,}1 & 3{,}5 & 2{,}7 & 3{,}0 & 2{,}5 & 2{,}1 & $-$10{,}9 & 6{,}3 & 6{,}9 & 2{,}4 & 3{,}8 & 2{,}7 \\
\hspace*{0.8em}Agricultura, ganadería, silvicultura y pesca & 7{,}5 & $-$2{,}3 & $-$3{,}3 & 2{,}6 & 5{,}2 & $-$9{,}2 & 13{,}8 & $-$1{,}2 & 4{,}8 & 4{,}5 & $-$3{,}5 & 4{,}2 & $-$2{,}8 & $-$2{,}0 & 7{,}0 & $-$16{,}9 & 4{,}7 & 15{,}0 & 11{,}9 \\
\hspace*{0.8em}Industria & 1{,}0 & $-$1{,}3 & $-$10{,}5 & 3{,}4 & $-$0{,}4 & $-$5{,}2 & $-$4{,}3 & 1{,}4 & 3{,}4 & 3{,}6 & 4{,}6 & 0{,}1 & 1{,}9 & $-$10{,}4 & 5{,}8 & 3{,}5 & $-$2{,}1 & 5{,}6 & 0{,}6 \\
\hspace*{1.6em}\textit{de la cual, manufacturas} & 0{,}5 & $-$2{,}8 & $-$11{,}4 & $-$0{,}7 & $-$1{,}2 & $-$5{,}7 & $-$1{,}3 & 2{,}2 & 5{,}1 & 1{,}7 & 6{,}8 & $-$1{,}1 & 0{,}6 & $-$14{,}1 & 13{,}9 & 6{,}5 & 0{,}4 & 3{,}7 & 2{,}2 \\
\hspace*{0.8em}Construcción & 1{,}2 & $-$1{,}7 & $-$8{,}5 & $-$15{,}3 & $-$13{,}2 & $-$9{,}7 & $-$10{,}1 & $-$0{,}3 & 6{,}0 & 3{,}5 & 1{,}7 & 3{,}0 & 4{,}7 & $-$14{,}7 & $-$1{,}0 & 8{,}9 & $-$0{,}3 & $-$0{,}9 & 6{,}0 \\
\hspace*{0.8em}Comercio, transporte y hostelería & 3{,}4 & 0{,}5 & $-$3{,}1 & 1{,}8 & 0{,}5 & $-$1{,}6 & $-$1{,}4 & 1{,}7 & 6{,}2 & 2{,}8 & 3{,}2 & 2{,}9 & 2{,}1 & $-$24{,}7 & 15{,}0 & 13{,}8 & 4{,}8 & 1{,}8 & 2{,}6 \\
\hspace*{0.8em}Información y comunicaciones & 4{,}0 & 1{,}9 & 0{,}1 & 2{,}8 & $-$0{,}7 & 1{,}5 & 2{,}8 & 6{,}9 & 6{,}1 & 4{,}2 & 8{,}5 & 4{,}1 & 4{,}2 & $-$4{,}8 & 8{,}4 & 14{,}6 & 7{,}2 & 6{,}8 & 4{,}9 \\
\hspace*{0.8em}Actividades profesionales y administrativas & 8{,}3 & 2{,}4 & $-$2{,}9 & $-$0{,}1 & 3{,}1 & $-$3{,}1 & $-$0{,}8 & 6{,}0 & 7{,}6 & 4{,}5 & 4{,}4 & 4{,}4 & 7{,}0 & $-$11{,}7 & 12{,}0 & 12{,}3 & 2{,}7 & 8{,}3 & 7{,}2 \\
\hspace*{0.8em}Administración pública, educación y sanidad & 2{,}9 & 3{,}2 & 1{,}8 & 1{,}6 & 2{,}0 & $-$2{,}0 & 0{,}2 & $-$0{,}9 & 0{,}9 & 1{,}1 & 2{,}2 & 1{,}4 & 1{,}4 & $-$1{,}5 & 1{,}9 & 1{,}5 & 2{,}5 & 3{,}7 & 1{,}3 \\
\bottomrule
\end{tabular}
\fuentecuadro{Eurostat (contabilidad nacional del INE).}
\end{bloquecuadro}
\clearpage

% ---- Página 4: cuadro 2 bis y gráficos 1 a 3 -------------------------------------
\setlength{\altograf}{6.2cm}\noindent
\begin{minipage}[t]{0.47\linewidth}\vspace{0pt}\setlength{\anchoetq}{7.4cm}
\begin{bloquecuadro}% Generado por generar_cuadros.py a partir de datos/cuadro2b-descomposicion.csv; no editar a mano.
\cuadro{2 bis}{Descomposición del crecimiento del PIB}{tasa media anual en \%; las filas sangradas suman la del PIB}
\setlength{\anchonum}{\dimexpr(\linewidth-\anchoetq)/3-2\tabcolsep\relax}
\begin{tabular}{@{}>{\raggedright\arraybackslash}p{\dimexpr\anchoetq-\tabcolsep\relax}*{3}{N}@{}}
\toprule
\rowcolor{cabecera} & \textbf{2007-2013} & \textbf{2013-2019} & \textbf{2019-2025} \\
\midrule
\textbf{PIB} & $-$1{,}33 & 2{,}59 & 1{,}59 \\
\hspace*{0.8em}Población & 0{,}49 & 0{,}18 & 0{,}79 \\
\hspace*{0.8em}Población en edad de trabajar (15-64) sobre población & $-$0{,}47 & $-$0{,}22 & $-$0{,}01 \\
\hspace*{0.8em}Tasa de ocupación (ocupados sobre población de 15-64) & $-$2{,}91 & 2{,}27 & 0{,}95 \\
\hspace*{0.8em}Horas trabajadas por ocupado & $-$0{,}07 & $-$0{,}10 & $-$0{,}46 \\
\hspace*{0.8em}PIB por hora trabajada & 1{,}63 & 0{,}46 & 0{,}33 \\
PIB por ocupado & 1{,}56 & 0{,}36 & $-$0{,}13 \\
PIB por habitante & $-$1{,}82 & 2{,}41 & 0{,}81 \\
PIB por persona en edad de trabajar & $-$1{,}35 & 2{,}63 & 0{,}82 \\
Horas trabajadas totales & $-$2{,}96 & 2{,}13 & 1{,}27 \\
\bottomrule
\end{tabular}
\fuentecuadro{Elaboración con datos de Eurostat (contabilidad nacional del INE: PIB en volumen, población, ocupados y horas; población de 15-64 años a 1 de enero). Diferencias logarítmicas.}
\end{bloquecuadro}
\end{minipage}\hfill
\grafico{1}{Contribuciones al crecimiento intertrimestral del PIB (pp)}{Eurostat (contabilidad nacional trimestral del INE, datos corregidos de estacionalidad y calendario).}{g01-contribuciones}
\par\vspace{10pt}\noindent
\grafico{2}{PIB real y PIB real por habitante (índice, 2007 = 100)}{Eurostat (contabilidad nacional del INE y de los países de la zona del euro).}{g02-pib-habitante}\hfill
\grafico{3}{Curva de Beveridge (medias anuales; 2026: primer semestre)}{Eurostat (estadística de vacantes, sin desestacionalizar; encuesta de población activa, tasa de paro desestacionalizada).}{g03-beveridge}
\clearpage

\setlength{\altograf}{5.6cm}
% ---- Página 5: cuadros 3 y 4 ----------------------------------------------------
\begin{bloquecuadro}% Generado por generar_cuadros.py a partir de datos/cuadro3.csv; no editar a mano.
\cuadro{3}{Mercado de trabajo, productividad, costes laborales y demografía}{en \%, salvo indicación}
\setlength{\anchonum}{\dimexpr(\linewidth-\anchoetq)/21-2\tabcolsep\relax}
\begin{tabular}{@{}>{\raggedright\arraybackslash}p{\dimexpr\anchoetq-\tabcolsep\relax}*{21}{N}@{}}
\toprule
\rowcolor{cabecera} & \textbf{2007} & \textbf{2008} & \textbf{2009} & \textbf{2010} & \textbf{2011} & \textbf{2012} & \textbf{2013} & \textbf{2014} & \textbf{2015} & \textbf{2016} & \textbf{2017} & \textbf{2018} & \textbf{2019} & \textbf{2020} & \textbf{2021} & \textbf{2022} & \textbf{2023} & \textbf{2024} & \textbf{2025} & \cabprev{2026} & \cabprev{2027} \\
\midrule
Empleo, personas (variación anual) & 3{,}2 & 0{,}1 & $-$6{,}3 & $-$1{,}8 & $-$2{,}4 & $-$4{,}3 & $-$2{,}2 & 1{,}0 & 2{,}9 & 2{,}1 & 2{,}6 & 2{,}2 & 2{,}7 & $-$4{,}4 & 2{,}6 & 3{,}8 & 3{,}1 & 3{,}0 & 2{,}6 & \prevcel{2{,}3} & \prevcel{1{,}4} \\
Horas trabajadas (variación anual) & 2{,}5 & 0{,}7 & $-$6{,}0 & $-$2{,}3 & $-$2{,}1 & $-$4{,}8 & $-$2{,}8 & 1{,}0 & 3{,}0 & 2{,}6 & 2{,}1 & 2{,}5 & 1{,}7 & $-$11{,}0 & 7{,}2 & 5{,}1 & 2{,}6 & 2{,}8 & 2{,}0 &  &  \\
Tasa de actividad (15-64 años) & 71{,}8 & 72{,}7 & 73{,}1 & 73{,}5 & 73{,}9 & 74{,}3 & 74{,}3 & 74{,}2 & 74{,}3 & 74{,}2 & 73{,}9 & 73{,}7 & 73{,}8 & 72{,}2 & 73{,}7 & 74{,}0 & 74{,}5 & 74{,}6 & 74{,}9 &  &  \\
Tasa de empleo (15-64 años) & 65{,}8 & 64{,}5 & 60{,}0 & 58{,}8 & 58{,}0 & 55{,}8 & 54{,}8 & 56{,}0 & 57{,}8 & 59{,}5 & 61{,}1 & 62{,}4 & 63{,}3 & 60{,}9 & 62{,}6 & 64{,}3 & 65{,}3 & 66{,}1 & 67{,}0 &  &  \\
\textbf{Tasa de paro (15-74 años)} & 8{,}2 & 11{,}3 & 17{,}9 & 19{,}9 & 21{,}4 & 24{,}8 & 26{,}1 & 24{,}5 & 22{,}1 & 19{,}6 & 17{,}2 & 15{,}3 & 14{,}1 & 15{,}5 & 14{,}9 & 13{,}0 & 12{,}2 & 11{,}4 & 10{,}5 & \prevcel{9{,}9} & \prevcel{9{,}6} \\
\hspace*{0.8em}Tasa de paro de la zona del euro &  &  & 9{,}7 & 10{,}3 & 10{,}3 & 11{,}5 & 12{,}2 & 11{,}7 & 11{,}0 & 10{,}1 & 9{,}1 & 8{,}2 & 7{,}6 & 8{,}0 & 7{,}8 & 6{,}8 & 6{,}6 & 6{,}4 & 6{,}4 &  &  \\
\hspace*{0.8em}Tasa de paro juvenil (15-24 años) & 18{,}1 & 24{,}5 & 37{,}7 & 41{,}5 & 46{,}2 & 52{,}9 & 55{,}5 & 53{,}2 & 48{,}3 & 44{,}4 & 38{,}6 & 34{,}3 & 32{,}5 & 38{,}3 & 35{,}0 & 29{,}7 & 28{,}7 & 26{,}5 & 24{,}9 &  &  \\
\hspace*{0.8em}Paro de larga duración (\% de los activos) &  &  & 4{,}3 & 7{,}3 & 8{,}9 & 11{,}0 & 13{,}0 & 12{,}9 & 11{,}4 & 9{,}5 & 7{,}7 & 6{,}4 & 5{,}3 & 5{,}0 & 6{,}2 & 5{,}1 & 4{,}3 & 3{,}8 & 3{,}4 &  &  \\
Asalariados temporales (\% de los asalariados) & 31{,}6 & 29{,}2 & 25{,}3 & 24{,}8 & 25{,}2 & 23{,}4 & 23{,}2 & 24{,}0 & 25{,}2 & 26{,}1 & 26{,}8 & 26{,}9 & 26{,}3 & 24{,}2 & 25{,}4 & 21{,}4 & 17{,}3 & 16{,}0 & 15{,}4 &  &  \\
Tasa de vacantes (\% de los puestos de trabajo) &  &  &  & 1{,}25 & 1{,}00 & 0{,}75 & 0{,}57 & 0{,}42 & 0{,}50 & 0{,}53 & 0{,}57 & 0{,}68 & 0{,}70 & 0{,}57 & 0{,}75 & 0{,}90 & 0{,}88 & 0{,}90 & 0{,}90 &  &  \\
Productividad por ocupado (variación anual) & 0{,}3 & 0{,}7 & 2{,}7 & 1{,}9 & 1{,}8 & 1{,}5 & 0{,}8 & 0{,}5 & 1{,}1 & 0{,}8 & 0{,}3 & 0{,}2 & $-$0{,}7 & $-$6{,}8 & 4{,}0 & 2{,}5 & $-$0{,}7 & 0{,}6 & 0{,}0 &  &  \\
Productividad por hora (variación anual) & 1{,}0 & 0{,}1 & 2{,}3 & 2{,}4 & 1{,}5 & 2{,}0 & 1{,}5 & 0{,}5 & 1{,}0 & 0{,}4 & 0{,}7 & $-$0{,}1 & 0{,}2 & 0{,}0 & $-$0{,}5 & 1{,}2 & $-$0{,}3 & 0{,}9 & 0{,}6 &  &  \\
Remuneración por asalariado (variación anual) & 4{,}0 & 6{,}8 & 4{,}0 & 1{,}2 & 0{,}0 & $-$1{,}1 & $-$0{,}1 & 0{,}2 & 1{,}2 & $-$0{,}1 & 0{,}7 & 1{,}7 & 3{,}4 & 0{,}8 & 4{,}8 & 4{,}7 & 5{,}1 & 3{,}7 & 4{,}4 &  &  \\
Costes laborales unitarios nominales (var. anual) & 3{,}7 & 6{,}1 & 1{,}3 & $-$0{,}7 & $-$1{,}8 & $-$2{,}6 & $-$0{,}9 & $-$0{,}3 & 0{,}0 & $-$0{,}8 & 0{,}5 & 1{,}5 & 4{,}1 & 8{,}2 & 0{,}8 & 2{,}1 & 5{,}9 & 3{,}0 & 4{,}4 &  &  \\
\hspace*{0.8em}Ídem, zona del euro & 1{,}4 & 3{,}9 & 4{,}4 & $-$0{,}5 & 0{,}5 & 2{,}1 & 1{,}2 & 0{,}6 & 0{,}3 & 0{,}8 & 0{,}7 & 2{,}0 & 1{,}9 & 4{,}4 & $-$0{,}4 & 3{,}2 & 6{,}5 & 4{,}2 & 3{,}3 &  &  \\
Población (variación anual) & 2{,}0 & 1{,}7 & 0{,}8 & 0{,}4 & 0{,}4 & 0{,}0 & $-$0{,}4 & $-$0{,}3 & $-$0{,}1 & 0{,}1 & 0{,}2 & 0{,}4 & 0{,}8 & 0{,}5 & 0{,}0 & 0{,}9 & 1{,}2 & 1{,}1 & 1{,}0 &  &  \\
Población extranjera (\% de la población) & 9{,}9 & 11{,}1 & 11{,}6 & 11{,}6 & 11{,}4 & 11{,}2 & 10{,}8 & 10{,}1 & 9{,}6 & 9{,}5 & 9{,}5 & 9{,}8 & 10{,}3 & 11{,}1 & 11{,}4 & 11{,}6 & 12{,}7 & 13{,}4 & 14{,}1 &  &  \\
Tasa de dependencia de los mayores (65+ / 15-64) & 24{,}0 & 23{,}8 & 24{,}1 & 24{,}6 & 25{,}2 & 25{,}7 & 26{,}3 & 27{,}2 & 27{,}8 & 28{,}2 & 28{,}6 & 29{,}0 & 29{,}2 & 29{,}5 & 29{,}7 & 30{,}2 & 30{,}4 & 30{,}8 & 31{,}2 &  &  \\
\bottomrule
\end{tabular}
\fuentecuadro{Eurostat (contabilidad nacional y encuesta de población activa del INE; estadística de vacantes; población a 1 de enero); previsiones: Comisión Europea (AMECO). La tasa de vacantes se da desde 2010 (salto de la serie entre 2009 y 2010).}
\end{bloquecuadro}
\vspace{10pt}
\begin{bloquecuadro}% Generado por generar_cuadros.py a partir de datos/cuadro4.csv; no editar a mano.
\setlength{\anchoetq}{7.6cm}
\cuadro{4}{Ahorro, inversión y financiación del sector privado}{en \%, salvo indicación}
\setlength{\anchonum}{\dimexpr(\linewidth-\anchoetq)/19-2\tabcolsep\relax}
\begin{tabular}{@{}>{\raggedright\arraybackslash}p{\dimexpr\anchoetq-\tabcolsep\relax}*{19}{N}@{}}
\toprule
\rowcolor{cabecera} & \textbf{2007} & \textbf{2008} & \textbf{2009} & \textbf{2010} & \textbf{2011} & \textbf{2012} & \textbf{2013} & \textbf{2014} & \textbf{2015} & \textbf{2016} & \textbf{2017} & \textbf{2018} & \textbf{2019} & \textbf{2020} & \textbf{2021} & \textbf{2022} & \textbf{2023} & \textbf{2024} & \textbf{2025} \\
\midrule
Tasa de ahorro de los hogares (\% de la renta disponible) & 6{,}0 & 8{,}1 & 11{,}7 & 10{,}0 & 10{,}6 & 7{,}0 & 8{,}5 & 7{,}0 & 7{,}9 & 7{,}7 & 6{,}3 & 6{,}1 & 8{,}6 & 17{,}6 & 14{,}3 & 9{,}0 & 11{,}4 & 12{,}8 & 11{,}9 \\
\hspace*{0.8em}Renta disponible real por habitante (variación anual) & 0{,}3 & 0{,}5 & 0{,}2 & $-$1{,}6 & $-$1{,}8 & $-$6{,}6 & $-$1{,}2 & 0{,}3 & 4{,}2 & 2{,}3 & 1{,}4 & 1{,}2 & 2{,}9 & $-$1{,}8 & 3{,}0 & $-$2{,}0 & 3{,}2 & 3{,}3 & 0{,}9 \\
Tasa de inversión de los hogares (\% de la renta disponible) & 15{,}8 & 13{,}9 & 10{,}8 & 9{,}8 & 7{,}9 & 6{,}6 & 4{,}9 & 4{,}8 & 4{,}7 & 4{,}7 & 5{,}2 & 5{,}5 & 5{,}6 & 5{,}3 & 6{,}4 & 7{,}5 & 6{,}8 & 7{,}0 & 7{,}4 \\
Tasa de inversión de las empresas no financieras (\% del VAB) & 30{,}2 & 25{,}7 & 18{,}8 & 19{,}3 & 19{,}6 & 21{,}9 & 22{,}8 & 24{,}3 & 23{,}9 & 24{,}9 & 25{,}6 & 27{,}0 & 27{,}6 & 27{,}4 & 25{,}8 & 23{,}9 & 24{,}3 & 24{,}8 & 24{,}4 \\
Participación de los beneficios de las empresas (\% del VAB) & 38{,}4 & 39{,}9 & 41{,}2 & 40{,}8 & 41{,}1 & 42{,}2 & 43{,}5 & 42{,}9 & 43{,}0 & 43{,}8 & 43{,}6 & 42{,}8 & 41{,}5 & 37{,}6 & 37{,}8 & 40{,}3 & 40{,}1 & 39{,}5 & 38{,}9 \\
Formación bruta de capital (\% del PIB) & 30{,}5 & 28{,}5 & 23{,}4 & 22{,}4 & 20{,}7 & 18{,}6 & 17{,}4 & 18{,}1 & 19{,}4 & 19{,}1 & 19{,}6 & 20{,}7 & 20{,}9 & 20{,}6 & 21{,}9 & 22{,}7 & 21{,}0 & 21{,}1 & 21{,}7 \\
Ahorro nacional bruto (\% del PIB) & 21{,}1 & 19{,}7 & 19{,}4 & 18{,}8 & 18{,}1 & 18{,}8 & 19{,}5 & 19{,}9 & 21{,}4 & 22{,}2 & 22{,}4 & 22{,}6 & 23{,}0 & 21{,}4 & 22{,}6 & 23{,}1 & 23{,}8 & 24{,}5 & 25{,}0 \\
Flujo de crédito al sector privado (\% del PIB) & 26{,}5 & 12{,}0 & $-$1{,}4 & 0{,}9 & $-$4{,}2 & $-$11{,}2 & $-$10{,}1 & $-$7{,}0 & $-$2{,}9 & 0{,}8 & 1{,}8 & $-$1{,}7 & 1{,}6 & 5{,}6 & 3{,}3 & $-$0{,}1 & 0{,}0 & 2{,}1 & 2{,}6 \\
Préstamos a hogares (var. interanual, diciembre) & 12{,}7 & 4{,}9 & $-$0{,}2 & 0{,}5 & $-$1{,}7 & $-$3{,}6 & $-$4{,}6 & $-$3{,}6 & $-$2{,}8 & $-$1{,}2 & $-$0{,}3 & 0{,}5 & $-$0{,}1 & $-$0{,}7 & 1{,}0 & 0{,}4 & $-$2{,}0 & 1{,}4 & 4{,}5 \\
Préstamos a empresas no financieras (var. interanual, dic.) & 19{,}7 & 6{,}9 & $-$3{,}3 & $-$1{,}1 & $-$4{,}1 & $-$10{,}9 & $-$11{,}5 & $-$7{,}1 & $-$0{,}9 & $-$0{,}6 & 0{,}4 & $-$2{,}7 & 0{,}1 & 9{,}1 & 1{,}6 & 0{,}6 & $-$3{,}7 & 0{,}3 & 4{,}2 \\
Tipo de los préstamos nuevos para vivienda & 5{,}09 & 5{,}67 & 3{,}25 & 2{,}51 & 3{,}34 & 3{,}27 & 2{,}99 & 2{,}93 & 2{,}23 & 1{,}94 & 1{,}91 & 1{,}94 & 1{,}98 & 1{,}73 & 1{,}48 & 1{,}97 & 3{,}78 & 3{,}40 & 2{,}73 \\
\hspace*{0.8em}Ídem, zona del euro & 4{,}93 & 5{,}30 & 4{,}01 & 3{,}44 & 3{,}74 & 3{,}47 & 3{,}07 & 2{,}80 & 2{,}26 & 1{,}97 & 1{,}86 & 1{,}82 & 1{,}63 & 1{,}39 & 1{,}31 & 2{,}08 & 3{,}67 & 3{,}70 & 3{,}30 \\
Tipo de los préstamos nuevos a empresas no financieras & 5{,}13 & 5{,}57 & 2{,}98 & 2{,}70 & 3{,}57 & 3{,}51 & 3{,}61 & 3{,}47 & 2{,}62 & 2{,}23 & 1{,}99 & 1{,}80 & 1{,}69 & 1{,}64 & 1{,}52 & 1{,}91 & 4{,}64 & 4{,}72 & 3{,}40 \\
Deuda de los hogares (\% del PIB) & 81{,}9 & 82{,}8 & 85{,}0 & 84{,}3 & 82{,}4 & 81{,}4 & 77{,}6 & 73{,}1 & 67{,}6 & 64{,}1 & 60{,}9 & 58{,}6 & 56{,}5 & 62{,}1 & 57{,}2 & 51{,}4 & 46{,}2 & 43{,}6 & 42{,}9 \\
Deuda de las empresas no financieras (\% del PIB) & 111{,}5 & 114{,}8 & 119{,}0 & 118{,}5 & 115{,}3 & 106{,}8 & 100{,}0 & 97{,}1 & 89{,}3 & 86{,}2 & 82{,}6 & 77{,}1 & 75{,}6 & 89{,}9 & 84{,}3 & 73{,}0 & 68{,}1 & 65{,}3 & 62{,}0 \\
Precio de la vivienda (variación anual) & 9{,}8 & $-$1{,}4 & $-$6{,}6 & $-$1{,}8 & $-$7{,}6 & $-$14{,}8 & $-$9{,}1 & 0{,}3 & 3{,}6 & 4{,}6 & 6{,}2 & 6{,}7 & 5{,}2 & 2{,}2 & 3{,}7 & 7{,}4 & 4{,}0 & 8{,}5 & 12{,}7 \\
\hspace*{0.8em}Precio de la vivienda deflactado (variación anual) & 6{,}2 & $-$4{,}8 & $-$5{,}7 & $-$3{,}7 & $-$9{,}7 & $-$16{,}5 & $-$10{,}0 & 0{,}2 & 3{,}7 & 4{,}4 & 4{,}6 & 5{,}2 & 4{,}1 & 2{,}2 & 1{,}4 & 0{,}7 & $-$1{,}0 & 5{,}2 & 9{,}7 \\
\bottomrule
\end{tabular}
\fuentecuadro{Eurostat (cuentas de los sectores institucionales y contabilidad nacional del INE; indicadores del procedimiento de desequilibrios macroeconómicos; índice de precios de la vivienda del INE) y Banco Central Europeo (préstamos y tipos de interés de las entidades de crédito; tipos: media anual, España salvo indicación).}
\end{bloquecuadro}
\clearpage

% ---- Página 6: gráficos 4 a 7 ---------------------------------------------------
\setlength{\altograf}{6.4cm}\noindent
\grafico{4}{Inflación: IPCA general y sin energía ni alimentos (tasa interanual, \%)}{Eurostat (índice de precios de consumo armonizado).}{g04-ipca}\hfill
\grafico{5}{Precio del petróleo Brent (media mensual, por barril)}{Administración de Información Energética de Estados Unidos (EIA); conversión a euros con el tipo de cambio de referencia del BCE.}{g05-brent}
\par\vspace{10pt}\noindent
\grafico{6}{Cuenta corriente y capacidad o necesidad de financiación (\% del PIB)}{Eurostat (balanza de pagos del Banco de España; contabilidad nacional del INE).}{g06-cuenta-corriente}\hfill
\grafico{7}{Posición de inversión internacional neta y deuda externa neta (\% del PIB)}{Eurostat (posición de inversión internacional del Banco de España).}{g07-piin}
\clearpage

% ---- Página 7: cuadro 5 y gráficos 8 y 9 ------------------------------------------
\begin{bloquecuadro}% Generado por generar_cuadros.py a partir de datos/cuadro5.csv; no editar a mano.
\cuadro{5}{Precios}{variación anual en \%}
\setlength{\anchonum}{\dimexpr(\linewidth-\anchoetq)/21-2\tabcolsep\relax}
\begin{tabular}{@{}>{\raggedright\arraybackslash}p{\dimexpr\anchoetq-\tabcolsep\relax}*{21}{N}@{}}
\toprule
\rowcolor{cabecera} & \textbf{2007} & \textbf{2008} & \textbf{2009} & \textbf{2010} & \textbf{2011} & \textbf{2012} & \textbf{2013} & \textbf{2014} & \textbf{2015} & \textbf{2016} & \textbf{2017} & \textbf{2018} & \textbf{2019} & \textbf{2020} & \textbf{2021} & \textbf{2022} & \textbf{2023} & \textbf{2024} & \textbf{2025} & \cabprev{2026} & \cabprev{2027} \\
\midrule
\textbf{IPCA: índice general} & 2{,}8 & 4{,}1 & $-$0{,}2 & 2{,}0 & 3{,}0 & 2{,}4 & 1{,}5 & $-$0{,}2 & $-$0{,}6 & $-$0{,}3 & 2{,}0 & 1{,}7 & 0{,}8 & $-$0{,}3 & 3{,}0 & 8{,}3 & 3{,}4 & 2{,}9 & 2{,}7 & \prevcel{3{,}0} & \prevcel{2{,}5} \\
\hspace*{0.8em}IPCA sin energía ni alimentos (subyacente) & 2{,}5 & 2{,}4 & 0{,}9 & 0{,}8 & 1{,}2 & 1{,}3 & 1{,}3 & $-$0{,}1 & 0{,}3 & 0{,}7 & 1{,}2 & 1{,}0 & 1{,}1 & 0{,}5 & 0{,}6 & 3{,}8 & 4{,}1 & 2{,}8 & 2{,}6 &  &  \\
\hspace*{0.8em}IPCA: energía & 1{,}7 & 11{,}9 & $-$9{,}0 & 12{,}5 & 15{,}7 & 8{,}8 & 0{,}0 & $-$0{,}8 & $-$9{,}0 & $-$8{,}6 & 8{,}0 & 6{,}1 & $-$1{,}2 & $-$9{,}6 & 21{,}1 & 27{,}6 & $-$16{,}1 & 0{,}9 & 3{,}1 &  &  \\
\hspace*{0.8em}IPCA: alimentos, bebidas alcohólicas y tabaco & 4{,}1 & 5{,}7 & 0{,}2 & 1{,}1 & 2{,}8 & 2{,}6 & 3{,}2 & $-$0{,}1 & 1{,}2 & 1{,}3 & 1{,}3 & 1{,}7 & 1{,}0 & 2{,}1 & 1{,}7 & 10{,}7 & 11{,}1 & 3{,}6 & 2{,}7 &  &  \\
IPCA de la zona del euro: índice general & 2{,}1 & 3{,}3 & 0{,}3 & 1{,}6 & 2{,}7 & 2{,}5 & 1{,}4 & 0{,}4 & 0{,}2 & 0{,}2 & 1{,}5 & 1{,}8 & 1{,}2 & 0{,}3 & 2{,}6 & 8{,}4 & 5{,}4 & 2{,}4 & 2{,}1 &  &  \\
\hspace*{0.8em}IPCA de la zona del euro: subyacente & 1{,}9 & 1{,}8 & 1{,}4 & 1{,}0 & 1{,}4 & 1{,}5 & 1{,}1 & 0{,}8 & 1{,}1 & 0{,}8 & 1{,}0 & 1{,}0 & 1{,}0 & 0{,}7 & 1{,}5 & 3{,}9 & 4{,}9 & 2{,}8 & 2{,}4 &  &  \\
Deflactor del PIB & 3{,}6 & 2{,}5 & 0{,}2 & 0{,}3 & $-$0{,}1 & $-$0{,}2 & 0{,}4 & $-$0{,}2 & 0{,}6 & 0{,}4 & 1{,}3 & 1{,}2 & 1{,}4 & 1{,}1 & 2{,}6 & 4{,}7 & 6{,}2 & 3{,}1 & 3{,}0 &  &  \\
Deflactor del consumo privado & 3{,}4 & 3{,}6 & $-$1{,}0 & 2{,}0 & 2{,}3 & 2{,}0 & 0{,}9 & 0{,}1 & $-$0{,}1 & 0{,}2 & 1{,}5 & 1{,}4 & 1{,}1 & 0{,}1 & 2{,}2 & 6{,}6 & 5{,}1 & 3{,}2 & 2{,}4 &  &  \\
Deflactor de las exportaciones & 2{,}5 & 2{,}7 & $-$2{,}7 & 3{,}5 & 4{,}3 & 2{,}3 & $-$0{,}7 & $-$1{,}7 & 0{,}5 & $-$1{,}2 & 2{,}5 & 1{,}8 & 0{,}5 & $-$0{,}9 & 6{,}9 & 14{,}6 & 1{,}5 & 1{,}2 & 1{,}0 &  &  \\
Deflactor de las importaciones & 1{,}6 & 4{,}4 & $-$7{,}0 & 6{,}5 & 8{,}2 & 3{,}4 & $-$1{,}9 & $-$1{,}0 & $-$0{,}1 & $-$1{,}6 & 3{,}1 & 2{,}5 & 0{,}7 & $-$3{,}0 & 7{,}6 & 22{,}4 & $-$4{,}5 & 0{,}0 & $-$0{,}6 &  &  \\
\bottomrule
\end{tabular}
\fuentecuadro{Eurostat (IPCA y contabilidad nacional del INE); previsión del IPCA: Comisión Europea (AMECO).}
\end{bloquecuadro}
\vspace{4pt}\noindent\setlength{\altograf}{4.6cm}
\grafico{8}{Tipo de cambio efectivo real y exportaciones (índice, 2007 = 100)}{Eurostat (tipos de cambio efectivos reales; contabilidad nacional del INE). Un aumento del tipo de cambio efectivo real indica una apreciación.}{g08-competitividad}\hfill
\grafico{9}{Exportaciones de bienes, suma de 12 meses (índice, 2024 = 100)}{Eurostat (Comext, estadística de comercio exterior de bienes). Valor en euros; suma de los últimos 12 meses en relación con su media de 2024.}{g09-exportaciones-eeuu}
\clearpage

% ---- Página 8: cuadro 6 y gráficos 10 y 11 ----------------------------------------
\begin{bloquecuadro}% Generado por generar_cuadros.py a partir de datos/cuadro6.csv; no editar a mano.
\setlength{\anchoetq}{7.6cm}
\cuadro{6}{Competitividad precio y no precio}{en \%, salvo indicación}
\setlength{\anchonum}{\dimexpr(\linewidth-\anchoetq)/19-2\tabcolsep\relax}
\begin{tabular}{@{}>{\raggedright\arraybackslash}p{\dimexpr\anchoetq-\tabcolsep\relax}*{19}{N}@{}}
\toprule
\rowcolor{cabecera} & \textbf{2007} & \textbf{2008} & \textbf{2009} & \textbf{2010} & \textbf{2011} & \textbf{2012} & \textbf{2013} & \textbf{2014} & \textbf{2015} & \textbf{2016} & \textbf{2017} & \textbf{2018} & \textbf{2019} & \textbf{2020} & \textbf{2021} & \textbf{2022} & \textbf{2023} & \textbf{2024} & \textbf{2025} \\
\midrule
Tipo de cambio efectivo real: IPC, 42 socios (2020 = 100) & 102{,}5 & 105{,}1 & 105{,}6 & 102{,}2 & 102{,}5 & 100{,}0 & 101{,}9 & 101{,}5 & 96{,}9 & 97{,}6 & 99{,}1 & 101{,}0 & 99{,}3 & 100{,}0 & 100{,}6 & 98{,}7 & 99{,}0 & 99{,}5 & 100{,}7 \\
Tipo de cambio efectivo real: CLU, zona del euro (2020 = 100) & 109{,}2 & 111{,}9 & 108{,}4 & 108{,}2 & 105{,}6 & 100{,}6 & 98{,}3 & 97{,}5 & 97{,}3 & 95{,}4 & 95{,}0 & 94{,}4 & 96{,}4 & 100{,}0 & 101{,}0 & 99{,}8 & 99{,}1 & 98{,}0 & 98{,}9 \\
Cuota en las exportaciones mundiales de bienes y servicios & 2{,}2 & 2{,}1 & 2{,}2 & 2{,}0 & 2{,}0 & 1{,}8 & 1{,}9 & 1{,}9 & 1{,}9 & 2{,}0 & 2{,}0 & 2{,}0 & 2{,}0 & 1{,}8 & 1{,}8 & 1{,}8 & 2{,}0 & 2{,}0 & 2{,}0 \\
\hspace*{0.8em}Bienes & 1{,}9 & 1{,}8 & 1{,}9 & 1{,}7 & 1{,}7 & 1{,}6 & 1{,}7 & 1{,}7 & 1{,}7 & 1{,}8 & 1{,}8 & 1{,}8 & 1{,}8 & 1{,}8 & 1{,}7 & 1{,}7 & 1{,}8 & 1{,}8 & 1{,}8 \\
\hspace*{0.8em}Servicios & 3{,}5 & 3{,}4 & 3{,}3 & 3{,}0 & 3{,}0 & 2{,}8 & 2{,}7 & 2{,}7 & 2{,}4 & 2{,}5 & 2{,}6 & 2{,}5 & 2{,}5 & 1{,}7 & 1{,}9 & 2{,}2 & 2{,}4 & 2{,}5 & 2{,}7 \\
Variación de la cuota de exportación en 5 años & $-$6{,}7 & $-$4{,}3 & 2{,}3 & $-$11{,}8 & $-$8{,}2 & $-$15{,}7 & $-$3{,}1 & $-$1{,}6 & 3{,}2 & 4{,}9 & 4{,}6 & 4{,}7 & $-$1{,}7 & $-$12{,}7 & $-$11{,}5 & $-$7{,}4 & 12{,}3 & 12{,}4 & 10{,}7 \\
Exportaciones de alta tecnología (\% de las exportaciones) & 4{,}2 & 4{,}2 & 4{,}8 & 4{,}8 & 4{,}8 & 5{,}0 & 5{,}4 & 5{,}2 & 5{,}5 & 5{,}8 & 5{,}7 & 5{,}5 & 5{,}8 & 6{,}0 & 7{,}1 & 8{,}5 &  &  &  \\
Gasto en I+D (\% del PIB) & 1{,}24 & 1{,}32 & 1{,}36 & 1{,}35 & 1{,}33 & 1{,}29 & 1{,}27 & 1{,}23 & 1{,}21 & 1{,}18 & 1{,}20 & 1{,}23 & 1{,}24 & 1{,}40 & 1{,}40 & 1{,}40 & 1{,}49 & 1{,}50 &  \\
\hspace*{0.8em}Ídem, Unión Europea & 1{,}79 & 1{,}87 & 1{,}96 & 1{,}96 & 2{,}00 & 2{,}06 & 2{,}08 & 2{,}09 & 2{,}10 & 2{,}10 & 2{,}14 & 2{,}17 & 2{,}21 & 2{,}28 & 2{,}24 & 2{,}22 & 2{,}25 & 2{,}26 &  \\
Apertura: (exportaciones + importaciones) / PIB & 57{,}6 & 55{,}8 & 46{,}8 & 52{,}7 & 58{,}5 & 60{,}6 & 61{,}7 & 63{,}4 & 63{,}6 & 63{,}2 & 66{,}2 & 67{,}0 & 66{,}4 & 59{,}4 & 66{,}5 & 78{,}5 & 71{,}9 & 70{,}1 & 69{,}1 \\
\bottomrule
\end{tabular}
\fuentecuadro{Eurostat (tipos de cambio efectivos reales; cuadro de indicadores del procedimiento de desequilibrios macroeconómicos; estadística de I+D del INE; contabilidad nacional). Un aumento del tipo de cambio efectivo real indica una apreciación real. Alta tecnología: hasta 2022.}
\end{bloquecuadro}
\vspace{8pt}\noindent\setlength{\altograf}{5.6cm}
\grafico{10}{Brecha de producción y variación del saldo estructural primario}{Comisión Europea (AMECO, previsiones de primavera de 2026). Saldo estructural primario y brecha, en \% del PIB potencial; 2026 y 2027 (círculos huecos): previsión.}{g10-orientacion-fiscal}\hfill
\grafico{11}{Deuda pública: tipo de interés implícito (i) y crecimiento del PIB nominal (g)}{Eurostat (cuentas de las administraciones públicas y contabilidad nacional). i: intereses del año entre la deuda del año anterior.}{g11-deuda-i-g}
\clearpage

% ---- Página 9: cuadro 7 y gráfico 12 ----------------------------------------------
\begin{bloquecuadro}% Generado por generar_cuadros.py a partir de datos/cuadro7.csv; no editar a mano.
\setlength{\anchoetq}{7.4cm}
\cuadro{7}{Política monetaria, crédito y sistema bancario}{en \%, salvo indicación}
\setlength{\anchonum}{\dimexpr(\linewidth-\anchoetq)/19-2\tabcolsep\relax}
\begin{tabular}{@{}>{\raggedright\arraybackslash}p{\dimexpr\anchoetq-\tabcolsep\relax}*{19}{N}@{}}
\toprule
\rowcolor{cabecera} & \textbf{2007} & \textbf{2008} & \textbf{2009} & \textbf{2010} & \textbf{2011} & \textbf{2012} & \textbf{2013} & \textbf{2014} & \textbf{2015} & \textbf{2016} & \textbf{2017} & \textbf{2018} & \textbf{2019} & \textbf{2020} & \textbf{2021} & \textbf{2022} & \textbf{2023} & \textbf{2024} & \textbf{2025} \\
\midrule
Facilidad de depósito del BCE (fin de año) & 3{,}00 & 2{,}00 & 0{,}25 & 0{,}25 & 0{,}25 & 0{,}00 & 0{,}00 & $-$0{,}20 & $-$0{,}30 & $-$0{,}40 & $-$0{,}40 & $-$0{,}40 & $-$0{,}50 & $-$0{,}50 & $-$0{,}50 & 2{,}00 & 4{,}00 & 3{,}00 & 2{,}00 \\
Operaciones principales de financiación (fin de año) & 4{,}00 & 2{,}50 & 1{,}00 & 1{,}00 & 1{,}00 & 0{,}75 & 0{,}25 & 0{,}05 & 0{,}05 & 0{,}00 & 0{,}00 & 0{,}00 & 0{,}00 & 0{,}00 & 0{,}00 & 2{,}50 & 4{,}50 & 3{,}15 & 2{,}15 \\
Euríbor a 12 meses (media anual) & 4{,}45 & 4{,}81 & 1{,}62 & 1{,}35 & 2{,}01 & 1{,}11 & 0{,}54 & 0{,}48 & 0{,}17 & $-$0{,}03 & $-$0{,}15 & $-$0{,}17 & $-$0{,}22 & $-$0{,}30 & $-$0{,}49 & 1{,}09 & 3{,}86 & 3{,}27 & 2{,}22 \\
M3 de la zona del euro (var. interanual, diciembre) & 11{,}5 & 7{,}7 & $-$0{,}3 & 1{,}2 & 1{,}6 & 3{,}6 & 1{,}0 & 3{,}7 & 4{,}7 & 5{,}0 & 4{,}6 & 4{,}2 & 4{,}8 & 12{,}3 & 6{,}9 & 3{,}8 & 0{,}0 & 3{,}6 & 2{,}8 \\
Tipo de los préstamos nuevos para vivienda & 5{,}09 & 5{,}67 & 3{,}25 & 2{,}51 & 3{,}34 & 3{,}27 & 2{,}99 & 2{,}93 & 2{,}23 & 1{,}94 & 1{,}91 & 1{,}94 & 1{,}98 & 1{,}73 & 1{,}48 & 1{,}97 & 3{,}78 & 3{,}40 & 2{,}73 \\
\hspace*{0.8em}Ídem, zona del euro & 4{,}93 & 5{,}30 & 4{,}01 & 3{,}44 & 3{,}74 & 3{,}47 & 3{,}07 & 2{,}80 & 2{,}26 & 1{,}97 & 1{,}86 & 1{,}82 & 1{,}63 & 1{,}39 & 1{,}31 & 2{,}08 & 3{,}67 & 3{,}70 & 3{,}30 \\
Tipo de los préstamos nuevos a empresas no financieras & 5{,}13 & 5{,}57 & 2{,}98 & 2{,}70 & 3{,}57 & 3{,}51 & 3{,}61 & 3{,}47 & 2{,}62 & 2{,}23 & 1{,}99 & 1{,}80 & 1{,}69 & 1{,}64 & 1{,}52 & 1{,}91 & 4{,}64 & 4{,}72 & 3{,}40 \\
Préstamos a hogares (var. interanual, diciembre) & 12{,}7 & 4{,}9 & $-$0{,}2 & 0{,}5 & $-$1{,}7 & $-$3{,}6 & $-$4{,}6 & $-$3{,}6 & $-$2{,}8 & $-$1{,}2 & $-$0{,}3 & 0{,}5 & $-$0{,}1 & $-$0{,}7 & 1{,}0 & 0{,}4 & $-$2{,}0 & 1{,}4 & 4{,}5 \\
Préstamos a empresas no financieras (var. interanual, dic.) & 19{,}7 & 6{,}9 & $-$3{,}3 & $-$1{,}1 & $-$4{,}1 & $-$10{,}9 & $-$11{,}5 & $-$7{,}1 & $-$0{,}9 & $-$0{,}6 & 0{,}4 & $-$2{,}7 & 0{,}1 & 9{,}1 & 1{,}6 & 0{,}6 & $-$3{,}7 & 0{,}3 & 4{,}2 \\
Deuda de los hogares (\% del PIB) & 81{,}9 & 82{,}8 & 85{,}0 & 84{,}3 & 82{,}4 & 81{,}4 & 77{,}6 & 73{,}1 & 67{,}6 & 64{,}1 & 60{,}9 & 58{,}6 & 56{,}5 & 62{,}1 & 57{,}2 & 51{,}4 & 46{,}2 & 43{,}6 & 42{,}9 \\
Deuda de las empresas no financieras (\% del PIB) & 111{,}5 & 114{,}8 & 119{,}0 & 118{,}5 & 115{,}3 & 106{,}8 & 100{,}0 & 97{,}1 & 89{,}3 & 86{,}2 & 82{,}6 & 77{,}1 & 75{,}6 & 89{,}9 & 84{,}3 & 73{,}0 & 68{,}1 & 65{,}3 & 62{,}0 \\
Bancos: capital ordinario de nivel 1, CET1 (\% APR) &  &  &  &  &  &  &  & 11{,}7 & 12{,}6 & 12{,}7 & 12{,}6 & 12{,}2 & 12{,}5 & 13{,}2 & 13{,}2 & 13{,}0 & 13{,}1 & 13{,}4 & 13{,}9 \\
\hspace*{0.8em}Ídem, Unión Europea &  &  &  &  &  &  &  & 15{,}8 & 13{,}6 & 14{,}1 & 15{,}0 & 14{,}8 & 15{,}0 & 15{,}8 & 15{,}8 & 15{,}5 & 16{,}1 & 16{,}4 & 16{,}4 \\
Bancos: capital total (\% de los APR) &  & 11{,}3 & 12{,}2 & 11{,}9 & 12{,}1 & 11{,}4 & 13{,}2 & 13{,}5 & 14{,}5 & 14{,}7 & 15{,}4 & 15{,}4 & 15{,}7 & 16{,}8 & 17{,}3 & 16{,}5 & 17{,}0 & 17{,}5 & 18{,}1 \\
Bancos: préstamos dudosos (\% de los préstamos) &  &  &  &  &  &  &  & 8{,}0 & 6{,}2 & 5{,}7 & 4{,}5 & 3{,}7 & 3{,}2 & 2{,}9 & 3{,}0 & 2{,}7 & 2{,}8 & 2{,}6 & 2{,}4 \\
\hspace*{0.8em}Ídem, Unión Europea &  &  &  &  &  &  &  & 6{,}7 & 6{,}8 & 5{,}3 & 4{,}1 & 3{,}2 & 2{,}8 & 2{,}6 & 2{,}1 & 1{,}9 & 1{,}9 & 2{,}0 & 2{,}0 \\
Bancos: rentabilidad de los recursos propios (ROE) &  & 12{,}7 & 9{,}2 & 8{,}8 & 0{,}2 & $-$25{,}6 & 6{,}0 & 5{,}7 & 6{,}6 & 5{,}1 & 7{,}1 & 8{,}3 & 6{,}8 & $-$3{,}6 & 10{,}5 & 10{,}1 & 12{,}0 & 13{,}9 & 14{,}2 \\
\hspace*{0.8em}Ídem, Unión Europea & 10{,}0 & $-$2{,}8 & 1{,}0 & 3{,}9 & $-$0{,}8 & $-$1{,}6 & 2{,}1 & 3{,}2 & 4{,}4 & 3{,}5 & 5{,}8 & 6{,}1 & 5{,}4 & 2{,}3 & 6{,}8 & 7{,}3 & 9{,}2 & 9{,}3 & 9{,}3 \\
\bottomrule
\end{tabular}
\fuentecuadro{Banco Central Europeo (tipos oficiales, tipos de interés de las entidades de crédito, balance de las entidades y datos bancarios consolidados) y Eurostat (deuda consolidada del sector privado). Tipos de los préstamos nuevos: media anual, España salvo indicación. APR: activos ponderados por riesgo.}
\end{bloquecuadro}
\vspace{6pt}\noindent\setlength{\altograf}{5.0cm}
\grafico{12}{Tipos de interés (\%, datos mensuales)}{Banco Central Europeo (tipos oficiales, Euríbor y tipos de interés de las entidades de crédito) y Eurostat (rendimientos de la deuda pública a 10 años).}{g12-tipos}\hfill
\begin{minipage}[t]{0.47\linewidth}\vspace{30pt}\centering
\textbf{FIN DEL CUADERNILLO}
\end{minipage}

\end{document}
